\documentclass[11pt]{article}
\usepackage{amsmath,amssymb,amsthm,mathtools,geometry,booktabs,enumitem}
\geometry{margin=1in}
\title{Step 117: Residual Coefficient-Map Instantiation for $\Xi^{\rm BC}$}
\author{RH Membrane Construction Notes}
\date{May 14, 2026}

\newtheorem{theorem}{Theorem}
\newtheorem{lemma}{Lemma}
\newtheorem{definition}{Definition}
\newtheorem{proposition}{Proposition}
\newtheorem{warning}{Warning}

\begin{document}
\maketitle

\section{Purpose}
The previous step reduced the remaining boundary-to-co-Poisson failure to the residual
\[
  \Xi^{\rm BC}_{\ell,a}=\mathfrak B_{\ell,a}^*\Pi_{Y_a}\mathfrak B_{\ell,a}\succeq0.
\]
The present step instantiates the finite residual coefficient map
\[
  R_N:Y_{R,N}\longrightarrow \mathbb C^{I_N}
\]
which is required before a Hecke/Dirichlet source frame can absorb this residual.

The main point is that the visibility constant is not a column norm lower bound on a fixed map. It is a spanning residual of a declared, constructible coefficient family. This is the exact finite analogue of the adequacy residual $\Xi$.

\section{Residual sector}
Let $Y_R$ denote the completed residual sector generated by the positive part of $\Xi^{\rm BC}$. A finite residual window is a pair
\[
  Y_{R,N},\qquad M_{R,N}:Y_{R,N}\to H_N^{\rm resp},
\]
where $H_N^{\rm resp}$ is the finite response core and
\[
  G_{R,N}=M_{R,N}^*M_{R,N}
\]
is the residual Gram form.

\begin{definition}[Declared coefficient dictionary]
A residual coefficient dictionary is a synthesis map
\[
  D_N:\mathbb C^{I_N}\to H_N^{\rm resp}
\]
whose atoms are declared before looking at the final residual test. In the intended construction, $D_N$ is a Burnol/co-Poisson or Burnol-to-Dirichlet hybrid dictionary with support, parity, Mellin-vanishing, semilocal-transport, and all-six records.
\end{definition}

Let $P_N$ be the orthogonal projection onto $\operatorname{Ran}D_N$, and define
\[
  R_N=D_N^\dagger M_{R,N},
  \qquad
  E_{{\rm coef},N}=(I-P_N)M_{R,N}.
\]

\section{Finite visibility identity}
\begin{theorem}[Residual coefficient visibility identity]
Let $H_N=D_N^*D_N$. Then
\[
  R_N^*H_NR_N
  =M_{R,N}^*P_NM_{R,N}
  =G_{R,N}-E_{{\rm coef},N}^*E_{{\rm coef},N}.
\]
Consequently the worst-direction residual visibility is
\[
  c_{R,N}
  =\lambda_{\min}\left(G_{R,N}^{-1/2}R_N^*H_NR_NG_{R,N}^{-1/2}\right)
  =1-\left\|E_{{\rm coef},N}G_{R,N}^{-1/2}\right\|^2.
\]
\end{theorem}

\begin{proof}
Because $R_N=D_N^\dagger M_{R,N}$ and $H_N=D_N^*D_N$, the form $R_N^*H_NR_N$ is exactly the pullback of the orthogonal projection onto $\operatorname{Ran}D_N$:
\[
  R_N^*H_NR_N=M_{R,N}^*P_NM_{R,N}.
\]
Since $I=P_N+(I-P_N)$ orthogonally,
\[
  M_{R,N}^*M_{R,N}=M_{R,N}^*P_NM_{R,N}+M_{R,N}^*(I-P_N)M_{R,N}.
\]
The second term is $E_{{\rm coef},N}^*E_{{\rm coef},N}$, proving the identity. The formula for $c_{R,N}$ follows by normalizing with $G_{R,N}^{-1/2}$.
\end{proof}

\begin{warning}[Kernel obstruction]
If $\ker R_N\neq0$, then no source frame acting through $R_N$ can charge all directions in $Y_{R,N}$. This is an adequacy failure, not a shortage of arithmetic strength.
\end{warning}

\section{Residual source-frame implication}
Let
\[
  G_{\mathcal X,N}=\sum_{\omega\in\mathcal X_N}\lambda_{\omega,N}
  Q_{\omega,N}^*\Theta_{\omega,N}^{-1}Q_{\omega,N}
\]
be the Hecke/Dirichlet source Gram on coefficient space.

\begin{theorem}[Residual matrix lower frame]
If
\[
  G_{\mathcal X,N}\succeq \gamma_N H_N
\]
and
\[
  R_N^*H_NR_N\succeq c_{R,N}G_{R,N},
\]
then
\[
  R_N^*G_{\mathcal X,N}R_N\succeq \gamma_Nc_{R,N}G_{R,N}.
\]
Thus the finite residual source strength is
\[
  \Lambda_{R,N}=\gamma_Nc_{R,N}.
\]
\end{theorem}

The completed route requires
\[
  \gamma_Nc_{R,N}\to\infty
\]
together with finite-window tail/exhaustivity:
\[
  G_R\preceq \Pi_N^*G_{R,N}\Pi_N+T_{R,N},
  \qquad
  \operatorname{tr}T_{R,N}\to0.
\]

\section{Literature placement}
Burnol supplies the native carrier and zero-evaluator/co-Poisson decomposition: the co-Poisson subspace $P_a$ is the perpendicular complement of the zero-evaluator span $Y_a$ for $a<1$, and the evaluator systems become complete/minimal at the relevant thresholds. Heap--Soundararajan supply a scalar source-strength template through short Dirichlet polynomials and moment estimates. The semilocal Connes--Consani--Moscovici framework supplies the response geometry in which finite local Euler factors enter through Hardy--Titchmarsh local-factor multipliers. Conrey--Li remains a survival warning against collapsing the construction into a de Branges/RKHS positivity ansatz.

\section{Finite sanity model}
The finite toy model built:
\begin{itemize}[leftmargin=2em]
\item a discrete Sonin projection from time/Fourier cutoffs;
\item a shifted boundary block $B=P\tau_{\log 2}(I-P)$;
\item a zero-evaluator proxy projection $\Pi_Y$;
\item a finite residual sector from leading eigendirections of $B^*\Pi_YB$;
\item three coefficient dictionaries: short Dirichlet atoms, Burnol/co-Poisson-like atoms, and a hybrid dictionary.
\end{itemize}
These checks are algebra sanity checks only; they are not RH evidence.

The finite audit found that ordinary short-Dirichlet atoms had essentially zero worst-direction visibility in the toy residual sector. The Burnol/co-Poisson-like family improved visibility but still left a large worst-direction residual. The hybrid family improved mean visibility but did not close the worst-direction gate. This supports the methodological conclusion: the residual-sector coefficient map is a real $\Xi$ gate and cannot be replaced by average visibility.

\section{Conclusion}
Step 117 converts the residual-source route into the exact pair of obligations
\[
  G_{\mathcal X,N}\succeq\gamma_NH_N,
  \qquad
  R_N^*H_NR_N\succeq c_{R,N}G_{R,N},
\]
with
\[
  \gamma_Nc_{R,N}\to\infty.
\]
The Heap--Soundararajan side targets $\gamma_N$. The Burnol/co-Poisson residual coefficient map targets $c_{R,N}$. The next construction step is to replace the toy dictionary by a mathematically declared Burnol-to-Dirichlet hybrid atom family and attempt a genuine visibility theorem for the residual sector.

\end{document}
